% Why Are HPC Monitoring Signals Hard to Forecast?
% Working draft — arXiv format, venue not yet chosen.
%
% Build: make paper   (pdflatex + bibtex; latexmk is not installed on the
% development machine, so the Makefile chains the passes explicitly)
%
% Every number in this document must be traceable to a committed table under
% results/tables/ or data/derived/. Placeholders are marked \needsresult so
% they are impossible to miss in the compiled PDF.

\documentclass[11pt,a4paper]{article}

\usepackage[utf8]{inputenc}
\usepackage[T1]{fontenc}
\usepackage{lmodern}
\usepackage[margin=1in]{geometry}
\usepackage{amsmath,amssymb}
\usepackage{booktabs}
\usepackage{graphicx}
\graphicspath{{figures/}}
\usepackage{xcolor}
\usepackage[hidelinks]{hyperref}
\usepackage{natbib}

% Number and unit helpers.
%
% siunitx is deliberately NOT used. It is absent from a minimal TeX Live
% install, and a paper that will not compile without an extra package is a
% reproducibility problem for exactly the readers most likely to try. These do
% the little this document actually needs.
\newcommand{\num}[1]{#1}
\newcommand{\unit}[2]{#1\,#2}
\newcommand{\unitrange}[3]{#1--#2\,#3}

% Marks a value that has not been computed yet. Deliberately loud.
\newcommand{\needsresult}[1]{\textcolor{red}{\textbf{[#1]}}}
% Cites an entry in OPEN_ISSUES.md, keeping paper and register in sync.
\newcommand{\oi}[1]{\textsc{oi-#1}}

\title{Why Are HPC Monitoring Signals Hard to Forecast?\\
       \large Intrinsic Predictability Across Eight Telemetry Families
       of a Tier-0 Supercomputer}

\author{Mohsen Seyedkazemi Ardebili\\
        \small \texttt{mohsen.seyedkazemi@gmail.com}}

\date{Working draft --- \today}

\begin{document}
\maketitle

\begin{abstract}
Operational telemetry from high-performance computing systems resists
forecasting to a degree that is widely reported anecdotally and rarely
explained. The usual response is to increase model capacity. We test the
opposite position: that the difficulty is largely \emph{intrinsic to the
signals}, differs systematically between signal families, and can be quantified
without fitting any forecaster. It largely survives our own test, but only under
a criterion we did not start with, and establishing that required validating the
test itself against synthetic data with a known answer.

Using the public M100 ExaData release from CINECA's Marconi100 Tier-0
supercomputer, we characterise telemetry from eight heterogeneous plugin
families --- out-of-band sensors, operating-system counters, scheduler
aggregates, alarm states, facility power, room-level electrical measurement,
cooling-plant control registers and air handlers, and exogenous weather ---
collected from the same machine over the same 31 months. (A ninth family,
per-job submission records, is collected but is event-structured rather than a
time series, and is excluded; \oi{8}.) The families are held
fixed in system, workload and period while signal type varies.
For each series we compute model-free predictability measures spanning entropy,
periodicity, long memory, non-stationarity, and quantisation, and separately
measure the forecast skill achieved by a tiered suite ranging from naive
baselines to gradient-boosted trees under a strictly causal rolling-origin
protocol.

We report five findings. First, the monitoring infrastructure is itself
non-stationary: of \num{573} distinct (plugin, metric) pairs observed across the
campaign, only \num{104} are present in every month, and those are exactly the
out-of-band sensor set. Second, a large share of what is monitored carries no
forecastable information at all --- the fraction of series that never change
over a month ranges from \num{64.8}\% of alarm-state signals to \num{8.4}\% of
out-of-band sensors --- so any benchmark that averages over such series reports
the composition of its sample as much as the skill of its models. Third, model
capacity does not help: measured against a same-step naive baseline, median
skill by tier at a one-hour horizon runs \num{0.115} for naive forecasters,
\num{-0.064} for statistical ones, \num{-0.327} for neural networks and
\num{-0.613} for gradient-boosted trees, and beyond three hours every tier is
worse than repeating the last value.

Fourth, the measures do transfer to signal families they have never seen, but
in \emph{ordering} far more than in \emph{level}. A model trained on every
family except one ranks the held-out family's series correctly --- median
within-family Spearman \(+0.52\), positive in \emph{all} \num{15} families and
significant for \num{10} --- and beats the best family-blind
predictor of skill in \num{13} of \num{15} families, by a median \num{21}\%
reduction in error. Two families fall below that baseline, and only one
substantively: monotone error counters at \(-0.68\), with thermal marginally
negative at \(-0.01\). The substantive failure is the counters, which a drift
forecast predicts almost exactly and which no model-free measure anticipates.

Fifth, and separating the two readings of the capacity result, we estimate each
series' information-theoretic ceiling --- entropy rate by a Lempel-Ziv
estimator, then maximum predictability by Fano's inequality --- and compare it
to what persistence achieves on the same alphabet. The families split cleanly.
Error counters, system metadata, alarm states and setpoint registers carry
\num{0.03}--\num{0.34} bits per step and persistence already achieves \num{0.99}
or better: no model can improve them. The other nine families carry
\num{1}--\num{3} bits per step, of which persistence leaves \num{12} to \num{31}
percentage points of next-step accuracy unclaimed. So ``hard to forecast'' names
two different situations, and only the first is a property of the signals.
(The Lempel-Ziv estimator converges from above at these series lengths, so the
ceiling is an estimate rather than a bound and reported headroom is a lower
bound; we validate the direction against sources with known entropy rate.)

Two controls establish that this is a property of the data rather than of the
method, and a third shows why the choice of criterion decides the answer. The
identical pipeline recovers a synthetic portable law at \(R^2=+0.25\), matching
an oracle given the true generating parameter, and correctly finds nothing when
skill is permuted. Running the real protocol with arbitrary groups of matched
size transfers easily (median \(R^2\) \(+0.35\), every group positive,
\(p<0.001\) over \num{100} draws), so the difficulty is specific to crossing
family boundaries. Yet the raw \(R^2\) of the real data is \(-0.17\), which is
\emph{worse} than the \(-0.07\) of the permuted arm containing no relation at
all --- because \(R^2\) is scored against an oracle that already knows the
held-out family's mean skill, and its denominator collapses on families whose
series are alike. Practically, a model-free profile can rank which signals in an
unfamiliar family will be harder, and estimate how hard far better than nothing,
but not as well as knowing the family.
\end{abstract}

\section{Introduction}
\label{sec:intro}

Operators of large computing facilities increasingly collect holistic telemetry,
and increasingly want to \emph{predict} it rather than merely react to it:
thermal headroom, power capping, and proactive maintenance are all forecasting
problems. Reported accuracy on this class of signal is nonetheless poor relative
to standard time-series benchmarks, and the literature generally treats that as
a modelling deficiency to be fixed with a stronger model.

We invert the question. Instead of asking which model does best, we ask how much
predictability is present at all, measured independently of any forecaster, and
whether that measurement predicts the skill a forecaster subsequently achieves.
If it does, ``HPC telemetry is hard to forecast'' becomes a quantity with named
causes rather than folklore.

Our contributions are: (i) a quantitative census of signal heterogeneity across
a 31-month Tier-0 monitoring campaign, including the instability of the
monitoring system itself; (ii) a model-free predictability profile for each of
\num{1402} series drawn from eight plugin families; (iii) a strictly causal,
scale-free forecast benchmark over four model tiers on the same series; and
(iv) a test of whether the first predicts the second, cross-validated by holding
out entire signal families; and (v) a per-family estimate of the
information-theoretic ceiling on next-step accuracy, which separates families
that are at their limit from families whose predictability our models fail to
extract.

The answer is a qualified yes, and the qualification took two controls and a
change of criterion to establish. The measures transfer to families they have
never seen --- strongly in ordering, partially in level --- but under the
statistic we first reported, grouped cross-validated \(R^2\), the same data
looks like a flat null. That statistic scores against an oracle which already
knows the held-out family's mean, and it degenerates on families whose series
are alike. We report the correction as a finding in its own right
(\S\ref{sec:linking}), because grouped cross-validation summarised by \(R^2\) is
a common design and telemetry families are routinely homogeneous.

\section{Background}
\label{sec:background}

\subsection{Marconi100 and the ExaData campaign}

Marconi100 was a Tier-0 system at CINECA comprising 980 IBM Power AC922 nodes,
each with two 16-core POWER9 processors and four NVIDIA V100 accelerators. The
ExaData campaign published holistic telemetry for the machine from March 2020 to
September 2022, collected through the ExaMon framework and released as
Hive-partitioned Parquet, one archive per month.

We catalogue every partition of all 31 archives --- \num{10504} partitions
holding roughly \num{55} billion rows --- and extract the subset the study
requires. Because node-level series of one metric are highly redundant, and
because nodes sharing a rack, inlet air and workload are not independent
samples (\oi{9}), analysis runs on a fixed panel of \num{96} \emph{primary}
entities --- nodes, cooling devices or PLC panels --- whose membership is a
stable hash of the entity label. The hash sees only the label, so panel
membership is identical across every metric and month and cross-metric
comparisons land on the same physical machines. Families keyed by a secondary
tag yield proportionally more series than the cap suggests: Nagios identifies a
series by node \emph{and} check description and reaches \num{1484} distinct
series in a single metric-month, where the \num{96} bound applies to the node
dimension alone.

\subsection{The nine plugin families}
\label{sec:families}

The value of this dataset for the present question is that it holds the system
fixed while varying the signal. Table~\ref{tab:families} summarises the
families and the mechanism that generates each.

\begin{table}[t]
\centering
\caption{Telemetry families and their generating mechanisms. Nominal cadences
are as documented by the dataset; measured cadences differ substantially for
some families (\S\ref{sec:census}).}
\label{tab:families}
\small
\begin{tabular}{llll}
\toprule
Family & Measures & Nominal cadence & Generating mechanism \\
\midrule
IPMI      & BMC sensors: temperature, power, fans & \unit{20}{s} & thermal diffusion, quantised \\
Ganglia   & OS and GPU counters                   & \unitrange{20}{300}{s}, irregular & software instrumentation \\
SLURM     & cluster scheduler aggregates          & \unitrange{5}{10}{s} & user submission behaviour \\
Nagios    & service and hardware alarm states     & \unit{15}{min} & discrete state machine \\
Job table & per-job submission records            & event-driven & not a time series (\oi{8}) \\
Schneider & liquid-cooling PLC registers          & \unit{20}{s} & closed-loop control \\
Logics    & room power, PUE, chillers             & --- & facility aggregation \\
Vertiv    & air-handling units                    & \unit{10}{s} & closed-loop control \\
Weather   & outdoor conditions                    & \unit{10}{min} & exogenous \\
\bottomrule
\end{tabular}
\end{table}

\section{Related work}
\label{sec:related}

Four bodies of work bear on this study. Claims below were checked against the
source papers rather than their abstracts; the four we could obtain lawfully are
read and verified, and the remainder are cited only for definitions we
implement (\oi{10}).

\paragraph{The dataset.} M100 ExaData~\citep{borghesi2023m100} publishes 49.9\,TB
of holistic telemetry --- management, workload, facility and infrastructure ---
from Marconi100 across two and a half years. Two of its derived releases matter
here. The unsupervised anomaly-detection set aggregates \emph{IPMI} metrics into
15-minute periods with average, standard deviation, minimum and maximum per
window, plus a Nagios-derived label; our common analysis grid matches that
choice, though we apply it to all nine families rather than IPMI alone. The
thermal-hazard set reduces sampling from 20 seconds to 10 minutes, and the
authors state the trade explicitly: that rate ``will preserve room- and
node-level thermal transient but discard computing-component-level ones''. That
is a direct statement of what aggregation costs, and it applies to our results
as much as to theirs (\oi{4}).

\paragraph{What this telemetry is normally used for.} The dominant application
is anomaly \emph{detection} rather than forecasting.
\citet{borghesi2019anomaly} train per-node autoencoders on HPC monitoring data
and report detecting previously unseen anomalies at accuracies between 88\% and
96\%. Detection asks whether the present sample is unusual given a learned model
of normal behaviour; forecasting asks what the next sample will be. The former
can succeed on signals whose future values are not predictable at all, which is
one reason the difficulty we measure has attracted little attention.

\paragraph{Feature-based characterisation of time series.} Our method follows the
tradition of describing a series by a fixed feature vector rather than by a
fitted model. \citet{lubba2019catch22} reduce 4791 candidate features to 22 by
selecting for classification performance and minimal redundancy across 93
datasets, achieving a roughly 1000-fold reduction in computation for an average
accuracy loss of 7\%, and covering linear and non-linear autocorrelation,
successive differences, value distributions and outliers, and fluctuation
scaling. Our profile is smaller and chosen for interpretability by mechanism
rather than by classification performance, but the premise --- that a compact,
model-free description captures much of what matters about a series --- is
theirs.

\paragraph{Capacity and forecast accuracy.} Our finding that heavier models do
worse has a direct precedent outside HPC. \citet{zeng2023transformers} introduce
one-layer linear models and report that they outperform sophisticated
Transformer-based long-term forecasters ``in all cases, and often by a large
margin'' across nine real-world datasets, arguing that permutation-invariant
self-attention discards temporal ordering. Our result is the same shape in a
different domain and at the opposite end of the complexity range: on this
telemetry a persistence forecast beats gradient-boosted trees and neural
networks at every horizon tested. Where they compare a linear model against
Transformers, we find the useful baseline is simpler still.

\paragraph{What we did not verify.} MASE~\citep{hyndman2006another}, permutation
entropy~\citep{bandt2002permutation}, detrended fluctuation
analysis~\citep{peng1994mosaic} and the Theta
method~\citep{assimakopoulos2000theta} are cited for definitions we implement
and test against known-answer series. Their full texts are paywalled and were
not read; no claim is made here about their findings beyond the definitions, and
we attribute no result to them. Three of the four works we did read were
obtained as preprints rather than the versions of record we cite, which we note
as a limitation of the audit rather than of the works.

\section{Method}
\label{sec:method}

\subsection{Campaign census}
\label{sec:census-method}

Tar members are read header-only, so cataloguing all 31 archives costs seconds
and no disk. For each extracted partition we then compute a per-series census:
one row per (plugin, metric, entity, month) recording sample count, observed
cadence, cadence jitter, gap fraction, coverage, missingness, value range,
distinct-value count and quantisation step. The census is deliberately
descriptive --- it contains no predictability measure and no model --- so that
a later finding of between-family differences cannot be an artefact of having
looked first.

\subsection{The entity panel}

Node-level series of one metric are highly redundant --- nodes share racks,
inlet air, cooling zones, and workload --- and are therefore not independent
samples (\oi{9}). Analysis runs on a fixed panel of \num{96} entities whose
membership is a stable hash of the entity label under the project seed. Because
the hash sees only the label, panel membership is identical across every metric
and every month, so cross-metric comparisons are made on the same physical
machines rather than on a fresh random slice.

\subsection{A-priori signal taxonomy}

Labels are assigned from the plugin and metric \emph{name} alone, before any
predictability is measured, so that a finding of between-family differences
cannot be a relabelling of the result it is meant to explain.

\subsection{Model-free predictability measures}

Each series receives a profile of \num{21} measures grouped by the mechanism it
probes: entropy and complexity (permutation entropy~\citep{bandt2002permutation}, spectral entropy, sample
entropy~\citep{richman2000physiological}, Lempel--Ziv complexity, gzip
compression ratio); periodicity (dominant-period strength, daily seasonal
strength); memory (autocorrelation at lag one and one day, detrended fluctuation
exponent~\citep{peng1994mosaic}); stationarity (block-mean drift ratio, variance
ratio, Augmented Dickey--Fuller~\citep{dickey1979distribution} and
KPSS~\citep{kwiatkowski1992testing} \(p\)-values, and a mean-shift changepoint
rate obtained by binary segmentation~\citep{truong2020selective}); and
shape (jump fraction, repetition fraction, monotone fraction, reset rate,
excess kurtosis of differences). Four instrument covariates --- quantisation
step, interval jitter, gap fraction and coverage --- are carried separately so
their contribution can be isolated from signal structure.

Two of these were added in response to an earlier failure of this study's own
analysis rather than chosen a priori, and we flag the distinction. Monotone
fraction and reset rate target counter-like signals, which a drift forecaster
handles almost exactly and which the original profile could not distinguish from
genuinely hard signals. They help but do not repair the case: error counters
remain the worst-transferring family at a level-corrected \(R^2\) of
\(-1.05\), against \(-2.39\) raw. Each
returns \texttt{nan} rather than raising on degenerate input, because constant
series are the common case here rather than an edge case
(\S\ref{sec:constant}). Implementations are validated against series whose
answers are known by construction --- a ramp, a sine, white noise, a random walk
--- in the accompanying test suite.

No single measure suffices, and the reason is instructive. Permutation entropy
assigns \num{0.000} both to monotone counters, which have some \num{17850}
distinct values, and to alarm states, which have two: a monotone series has
exactly one ordering pattern and a step function almost never changes. Both are
easy to forecast, for opposite structural reasons, and entropy cannot
distinguish them (\oi{11}).

\subsection{Forecast protocol}

Skill is measured under rolling-origin evaluation with strict causality
asserted per split. Errors are scaled by MASE with a denominator computed from
the training window only; scaling by test-window statistics is the most common
silent leak in forecasting benchmarks and would compress exactly the
between-tier differences this study reads. The seasonal period is fixed per
series \emph{a priori} from its own cadence, with a non-seasonal robustness
check (\oi{5}).

MASE is undefined where the in-sample seasonal-naive error vanishes. This is
not an edge case here: a substantial share of control-plane series are constant
(\S\ref{sec:census}), and near-perfectly periodic windows yield denominators at
the level of floating-point noise. Such series are reported as undefined rather
than assigned a spuriously large score.

\subsection{Pipeline and cost}
\label{sec:pipeline}

The study is a thirteen-stage pipeline, each stage writing a committed artefact so
that any stage can be re-run without repeating the ones before it
(Table~\ref{tab:pipeline}). Bulk intermediates are sharded per partition and
skipped when present, which makes an interrupted run resumable rather than
repeatable.

Two engineering constraints shaped the design and are worth stating because they
bound what the study could attempt. The archives hold roughly \num{55} billion
rows against a machine with \num{62}\,GB of memory shared with other workloads,
so extraction is selective and concurrency is derived from a memory budget rather
than the core count. And every worker pins its BLAS and OpenMP pools to a single
thread: with pools left at their defaults, five worker processes on a 24-core
machine request 120 threads and the forecast sweep runs roughly \num{90} times
slower, with no error and no warning to indicate it. The full sweep takes minutes
once pinned.

\begin{table}[t]
\centering
\caption{Pipeline stages and their committed outputs.}
\label{tab:pipeline}
\small
\begin{tabular}{llp{2.9cm}}
\toprule
\# & Stage & Output \\
\midrule
1 & inventory & partition catalogue (all 31 archives) \\
2 & extract & selected partitions, disk-guarded \\
3 & census & per-series descriptive table \\
4 & taxonomy & a-priori labels + firing rule \\
5 & sample & stratified series sample \\
6 & features & \num{21}-measure profile per series \\
7 & forecast & skill per (series, model, horizon, origin) \\
8 & link & transfer, ablations, null control \\
9 & figures & figures + backing CSV per figure \\
10 & literature & screening table, bibliography \\
11 & direct control & recursive vs direct forecasting \\
12 & protocol control & synthetic arms with known ground truth \\
13 & ceiling & Fano bound per series (slowest stage, ${\sim}35$\,min) \\
\bottomrule
\end{tabular}
\end{table}

\begin{table}[t]
\centering
\caption{Protocol and model settings, fixed before results existed.}
\label{tab:hyper}
\small
\begin{tabular}{ll}
\toprule
Setting & Value \\
\midrule
Analysis grid & \num{15} min, common to all families \\
Entity panel & \num{96} primary entities, hashed \\
Horizons & 1, 4, 12, 96 steps (15\,min--24\,h) \\
Rolling origins & 4 per series \\
Training window & \num{1500} most-recent points \\
Minimum series length & \num{512} points \\
MASE denominators & same-step (primary), daily \\
Boosted trees & 8/15, 24/60, 48/200 lags/iters \\
Neural & 48--24 and 96--48 units, 40--60 iters \\
Ridge AR & order 24, \(\lambda=10^{-6}\) \\
Seed & \num{20230288} \\
\bottomrule
\end{tabular}
\end{table}

\subsection{Failure modes}
\label{sec:failures}

Degenerate input is the common case here, not an edge case, so the pipeline is
specified for it. A constant series yields an undefined MASE (the seasonal-naive
denominator is zero) and is excluded from skill statistics and counted; so is a
near-perfectly periodic one, whose denominator is at the level of floating-point
noise and would otherwise produce an arbitrary ratio of two rounding errors. Every
measure returns \texttt{nan} rather than raising on a series that is too short or
has no variation. A series whose value column never parses as numeric is recorded
as non-numeric and dropped. Where a forecaster fails to fit, the failure is
recorded against that (series, model, origin) and the sweep continues, so one
ill-conditioned fit cannot destroy a run.

\subsection{Linking analysis}

Skill is regressed on the predictability profile, cross-validated by holding out
an \emph{entire signal family}. A random split would let the model memorise
family-specific quirks and report a flattering \(R^2\); a held-out family forces
the relation to transfer to signal physics it has never seen, which is what a
portable difficulty law would require. Every \(R^2\) is reported against a
train-mean baseline, so ``the measures explain skill'' means better than knowing
nothing rather than better than zero. Three ablations exist to attack the
headline: dropping the instrument covariates, keeping only them, and removing
the tie-biased entropy. Marginal correlations use a cluster bootstrap over
metrics, because series of one metric across nodes are not independent
(\oi{9}).

\section{Results}
\label{sec:results}

\subsection{The monitoring system is not stationary}
\label{sec:census}

Across the 31-month campaign we catalogue \num{10504} metric partitions
comprising \num{573} distinct (plugin, metric) pairs. Only \num{104} --- all of
them IPMI --- are present in every month. Ganglia's metric count runs
\num{0}, \num{33}, \num{73}, \num{49}, \num{157}, \num{177} and is not
monotonic; SLURM's runs \num{20}, \num{2}, \num{0}, \num{2}, \num{36};
Schneider first appears in 21-03 and Vertiv and Weather in 21-04.

Figure~\ref{fig:coverage} shows the pattern. Any claim compared across months
therefore confounds signal behaviour with changes to the instrument, and we
restrict cross-month statements accordingly (\oi{2}).

\begin{figure}[t]
\centering
\includegraphics[width=\columnwidth]{coverage.pdf}
\caption{Metrics collected per plugin per month. Grey marks a plugin absent that
month; the colour scale is logarithmic.}
\label{fig:coverage}
\end{figure}

\subsection{Much of what is monitored never changes}
\label{sec:constant}

Across \num{159506} series the fraction constant over a whole month varies by
more than a factor of seven between families: \num{64.8}\% for Nagios,
\num{47.0}\% Schneider, \num{42.9}\% Logics and Ganglia, \num{33.1}\% Vertiv,
against \num{10.5}\% SLURM, \num{8.4}\% IPMI and \num{0}\% Weather. MASE is
undefined for such series --- the seasonal-naive denominator is zero --- and
they are excluded from skill statistics and counted rather than silently
dropped. Figure~\ref{fig:constant} gives the breakdown.

\begin{figure}[t]
\centering
\includegraphics[width=\columnwidth]{constant_share.pdf}
\caption{Share of series that never change over a month, by family.}
\label{fig:constant}
\end{figure}

\subsection{Regularity and interruption are different properties}
\label{sec:sampling}

Sampling irregularity is routinely summarised by the coefficient of variation of
the inter-sample interval. On this data that measure is misleading: it reports
IPMI as the least regular family, contradicting the dataset documentation. IPMI
in fact has \emph{zero} series with a jittery interval and \num{100}\% with a
gap exceeding a hundred times the median --- its sampling is exact and
universally interrupted, and the rare enormous gaps dominate the coefficient of
variation. Ganglia is the only family whose interval genuinely wanders
(\num{81}\% of series), which is what the documentation reports. We therefore
measure jitter and interruption separately (\oi{1}); Figure~\ref{fig:sampling}
shows that the two are nearly orthogonal.

\begin{figure}[t]
\centering
\includegraphics[width=\columnwidth]{sampling_irregularity.pdf}
\caption{Interval jitter and outage incidence are different properties. IPMI has
no jittery series and no uninterrupted ones.}
\label{fig:sampling}
\end{figure}

\subsection{Forecast skill and model capacity}
\label{sec:skill}

Figure~\ref{fig:tier} shows median skill by tier and horizon on \num{1082}
series, pooling every model within a tier. Capacity does not buy skill: the
naive tier leads at every horizon and the best single model throughout is
persistence.

Pooling penalises a tier for containing a deliberately weak member, so we also
compare each tier's \emph{best} model per series, paired across series. That
comparison is both more favourable to the sophisticated tiers and more
conclusive against them. At a one-hour horizon the per-series best is
\num{0.945} (IQR \([0.70, 1.00]\)) for naive, \num{0.861} for statistical,
\num{0.747} for boosted trees and \num{0.643} for neural networks; a paired
Wilcoxon test of naive against each other tier rejects equality at every horizon
and by a wide margin (\(p < 10^{-20}\) against every tier at one hour, and
\(p < 10^{-90}\) against the neural tier). The ordering is not an artefact of
pooling. At a 24-hour horizon even the best naive model is negative
(\num{-0.284}), so the claim that no tier beats a same-step naive forecast holds
at a day but not at three hours, where the per-series best naive model still
scores \num{0.880}. One qualification remains: the ordering depends on the baseline. Against a
daily-seasonal denominator the statistical tier leads instead, and boosted trees
overtake naive at the shortest horizon (\oi{5}).

\subsection{Control: is this an artefact of recursive forecasting?}
\label{sec:direct}

Every learned model above forecasts recursively --- predict a step, append it,
predict again --- so error compounds, and the collapse of the learned tiers at
long horizons might be that compounding rather than anything about the signals.
This is the most plausible alternative explanation for the result, so we test it
rather than caveat it. We fit the same gradient-boosted learner two ways on the
same series, origins and horizons: recursively, and \emph{directly}, trained to
map a lag window onto the value \(h\) steps ahead in a single prediction.

\begin{table}[t]
\centering
\caption{Median skill (\(1-\)MASE, same-step denominator) for the same learner
forecast recursively and directly, against persistence, on \num{341} matched
(series, origin) windows.}
\label{tab:direct}
\small
\begin{tabular}{lrrr}
\toprule
Horizon & Persistence & GBT recursive & GBT direct \\
\midrule
15\,min & \num{0.767} & \num{0.320} & \num{0.392} \\
1\,h    & \num{0.499} & \num{-0.223} & \num{0.114} \\
3\,h    & \num{-0.023} & \num{-1.029} & \num{-0.259} \\
24\,h   & \num{-2.758} & \num{-3.692} & \num{-2.838} \\
\bottomrule
\end{tabular}
\end{table}

Recursion does cost a great deal --- forecasting directly recovers between
\num{0.07} and \num{0.85} of skill depending on horizon, most at one and three
hours --- so the caveat was warranted. But the conclusion is unaffected: the
direct model still loses to persistence at every horizon tested
(Table~\ref{tab:direct}). The capacity result is a property of the problem, not
of the recursion.

\begin{figure}[t]
\centering
\includegraphics[width=\columnwidth]{skill_by_tier.pdf}
\caption{Median skill by model tier. Zero is a same-step naive forecast.}
\label{fig:tier}
\end{figure}

\subsection{Two regimes: signal-limited and model-limited families}
\label{sec:ceiling}

The capacity result above says the models we tried do not help. It does not say
whether anything could. Those are different claims, and separating them needs a
quantity no forecaster can beat.

We estimate one. Each series is symbolised at its own resolution --- its
distinct values where there are at most \num{16}, equal-frequency quantiles
otherwise --- its entropy rate \(S\) is estimated with the Lempel-Ziv estimator
\(S = n\log_2 n / \sum_i \Lambda_i\), where \(\Lambda_i\) is the length of the
shortest substring at position \(i\) not seen before it, and the maximum
predictability \(\Pi_{\max}\) follows from the limiting case of Fano's
inequality
\[
S = -[\Pi\log_2\Pi + (1-\Pi)\log_2(1-\Pi)] + (1-\Pi)\log_2(N-1),
\]
after \citet{song2010limits}. Against it we place the accuracy persistence
actually achieves on the \emph{same} alphabet --- persistence because
\S\ref{sec:skill} finds it the strongest forecaster here, so the comparison
is not against a straw man.

\textbf{This is an estimate, not a bound, and we validated the direction of its
error rather than assuming it.} On Markov chains whose entropy rate is known in
closed form, the Lempel-Ziv estimator converges from \emph{above} at the series
lengths available: at \(p(\text{stay})=0.99\) the true rate is \num{0.076} bits
and the estimate at \(n=2000\) is \num{0.127}, closing only near
\(n=16000\) --- far beyond the \(\approx\num{3000}\) samples a month of
15-minute data provides. The ceiling is therefore \emph{under}-reported by
\(\approx\num{0.01}\), and \num{248} of \num{1402} series (\num{18}\%) exceed
it. That bias direction is what makes the estimate usable: reported headroom is
a \emph{lower} bound on true headroom, so "these models leave predictability
unclaimed" is the conservative reading, while "this family is at its limit"
cannot be established from a small positive gap and we assert it only where the
gap is zero or negative.

\begin{table}[t]
\centering
\caption{Achieved next-symbol accuracy against the estimated information
ceiling, by family. Headroom is a lower bound. Four families are at their
limit; nine are not.}
\label{tab:ceiling}
\small
\begin{tabular}{lrrrr}
\toprule
Regime (families) & \(S\) & ceiling & achieved & headroom \\
\midrule
signal-limited (4)  & 0.03--0.34 & 0.97--1.00 & 0.99--1.00 & \(\approx 0\) \\
near-limit (4)      & 0.63--0.93 & 0.89--0.93 & 0.81--0.90 & 0.02--0.08 \\
model-limited (9)   & 1.00--2.97 & 0.50--0.88 & 0.18--0.69 & 0.12--0.31 \\
\bottomrule
\end{tabular}
\end{table}

The families split into two regimes (Table~\ref{tab:ceiling};
\texttt{results/tables/ceiling\_gap.csv}). Error counters, system metadata,
alarm states and setpoint registers carry \num{0.03}--\num{0.34} bits per step,
and persistence already achieves \num{0.99} or better --- at or above the
estimated ceiling. Nothing can improve them, and a benchmark that includes them
is measuring its own sample composition, which is the same point
\S\ref{sec:constant} makes from a different direction.

The other nine are not at any limit. Power, network, electrical, thermal, fan
flow, utilisation and humidity signals carry \num{1}--\num{3} bits per step, of
which persistence leaves \num{12} to \num{31} percentage points of next-step
accuracy unclaimed.

\textbf{This qualifies \S\ref{sec:skill} rather than contradicting it.} Both
hold: the model classes we tried --- boosted trees and lag-window networks,
forecasting recursively --- do worse than persistence, \emph{and} substantial
extractable structure demonstrably remains in the continuous physical families.
The defensible conclusion is therefore not that these signals are
unforecastable, but that the model classes tried are the wrong ones for them.
That is a sharper claim than a null, and it points somewhere.

One caveat bounds what the comparison licenses. The ceiling concerns
\emph{one-step} next-symbol accuracy on a \num{16}-level alphabet; our skill
numbers are MASE at a \emph{four-step} horizon on continuous values. Headroom
establishes that the information is present, not that a forecaster can convert
it into MASE.

\subsection{Does the protocol work? A synthetic positive control}
\label{sec:protocol}

Before reporting a negative transfer result we must establish that the protocol
can report a positive one. Otherwise ``we found no portable law'' is
indistinguishable from ``our procedure cannot find one''. We therefore generate
synthetic families whose ground truth is known by construction and run them
through the \emph{identical} pipeline --- the same measure functions, the same
forecast harness, the same leave-one-family-out routine --- varying only how the
series were produced (Table~\ref{tab:protocol}).

\begin{table}[t]
\centering
\caption{Synthetic protocol control, \num{600} series over \num{15} families
per arm, against the real result. The protocol recovers a genuine portable law
(\texttt{portable}); a disjoint-band world destroys the ordering, and M100 does
not look like it.}
\label{tab:protocol}
\small
\begin{tabular}{lrrr}
\toprule
Arm & mean \(R^2\) & fam.\ \(>0\) & med.\ Spearman \\
\midrule
\texttt{portable}     & \num{+0.248} & 13/15 & \num{+0.584} \\
\texttt{portable\_oracle} & \num{+0.233} & 13/15 & \num{+0.576} \\
\texttt{disjoint}     & \num{-0.096} & 2/15  & \num{-0.050} \\
\texttt{no\_law}       & \num{-0.070} & 6/15  & \num{+0.100} \\
\midrule
\textbf{M100 (real)}  & \num{-0.189} & 9/15  & \num{+0.495} \\
\bottomrule
\end{tabular}
\end{table}

Three things follow. First, the protocol works: where a portable law exists it
is recovered at \(R^2 = +0.25\) with \num{13} of \num{15} families positive,
so M100's negative \(R^2\) is a statement about M100 rather than about the
method. Second, supplying the \emph{true} generating parameter as a feature
(\texttt{portable\_oracle}) does not beat the measured features
(\num{+0.233} against \num{+0.248}), which bounds how much of the remaining gap
could be closed by better features alone.

Third, and least comfortable for the obvious explanation: M100 does \emph{not}
resemble the \texttt{disjoint} arm. If each family occupied its own
non-overlapping band of a shared law, predicting a held-out family would require
extrapolation and the ordering would fail with it --- and it does, to a median
Spearman of \(-0.05\). M100's ordering instead sits at \(+0.50\), close to the
genuinely portable arm's \(+0.58\). Whatever prevents the level from
transferring, it is not simply that families occupy disjoint parameter ranges.
We report this as an open mechanism rather than adopt the explanation the arm
was built to test.

\paragraph{What blocks the level.} The synthetic \texttt{disjoint} arm rules
out the crudest version of the extrapolation story, but not the idea itself. For
each held-out family we measure how far its feature distribution sits outside
the training families' support. That distance predicts the level error strongly
(\(r=+0.78\) against the absolute level bias) and the raw \(R^2\) nearly as
strongly (\(-0.66\)), while barely touching the ordering (\(-0.22\)). Error
counters sit furthest outside the training support of any family and are also
the worst transfer case on every metric.

The mechanism is therefore extrapolation, but selectively: predicting a family
whose measures lie outside anything seen in training biases the \emph{level} of
the prediction without destroying the \emph{ranking}. The synthetic arm failed
to reproduce this because forcing families into non-overlapping bands removes
the shared structure the ranking depends on, which is a stronger intervention
than the real data applies.

\paragraph{How much local labelling repairs it.} If the level failure is
extrapolation, a handful of labelled series from the unfamiliar family should
fix it, and the families needing most extrapolation should benefit most. We test
this directly: for each held-out family we add \(k\) of its own series to the
training set and re-fit, averaging over draws
(Table~\ref{tab:fewshot}).

\begin{table}[t]
\centering
\caption{Few-shot calibration. Adding \(k\) labelled series from the held-out
family before re-fitting. Twenty examples turn the mean \(R^2\) positive.}
\label{tab:fewshot}
\small
\begin{tabular}{rrrr}
\toprule
\(k\) & mean \(R^2\) & median \(R^2\) & families \(>0\) \\
\midrule
0  & \num{-0.212} & \num{+0.057} & 60\% \\
1  & \num{-0.170} & \num{+0.054} & 61\% \\
2  & \num{-0.142} & \num{+0.053} & 64\% \\
5  & \num{-0.090} & \num{+0.054} & 68\% \\
10 & \num{-0.015} & \num{+0.056} & 70\% \\
20 & \num{+0.061} & \num{+0.097} & 72\% \\
\bottomrule
\end{tabular}
\end{table}

Twenty labelled series --- a few hours of one operator's attention --- move the
mean \(R^2\) from \(-0.21\) to \(+0.06\) and raise the share of families
above zero from \num{60}\% to \num{72}\%. The gain is concentrated exactly
where the mechanism predicts. Error counters, furthest outside the training
support, improve from \(-2.39\) to \(-0.03\); alarm states from \(-2.41\) to
\(-1.15\). Families that already transferred barely move: CPU utilisation goes
\(+0.18\) to \(+0.12\) and the scheduler family \(+0.47\) to \(+0.51\).

This is the practical form of the result. The profile ranks candidate signals in
an unfamiliar family with no local data at all; calibrating the absolute
difficulty needs on the order of twenty labelled examples from that family, and
needs them most precisely where the ranking was already usable but the level was
not.

\paragraph{An information-theoretic reference point, and what it is not.} A
natural question is how far the achieved skill sits from an information-theoretic
limit. We compute one for every series: the entropy rate of the symbolised
signal, converted through Fano's inequality into an upper bound on the fraction
of next-bin transitions any predictor could get right. Computed on a 16-symbol alphabet, the median across
\num{1402} series is \num{0.854}, and it separates families sharply --- network
traffic at \num{0.575} and electrical measurements at \num{0.638} against alarm
states and setpoint registers at \num{0.996}.

It is \emph{not} a bound on our skill metric, and we report that rather than
let the framing imply otherwise. \num{77.5}\% of series achieve a skill above
their own ``ceiling''. The two quantities answer different questions: Fano
bounds next-\emph{bin} accuracy on a symbolised sequence, while skill is a
scale-free error ratio on continuous values, and a signal that is easy to
extrapolate but hard to bin exactly can exceed the first while doing well on the
second. Anyone reading the number as a limit on achievable error would be
misled.

What it is instead is the single strongest ordinal predictor of difficulty we
found: Spearman \(+0.69\) with achieved skill across series and \(+0.73\)
across family medians, ahead of the compression ratio and every entropy measure.
That is consistent with the paper's central result --- these measures rank well
and calibrate badly --- and it extends it, since the ceiling is derived from
information theory rather than fitted to anything.

\paragraph{A scale-free view.} Because \(R^2\) is dominated by that level
error, we also report mean absolute error against the same baseline. On that
metric the model beats predicting the training mean for \num{13} of \num{15}
families (mean skill \num{+0.20}), against \num{9} of \num{15} on \(R^2\).
Both are correct; they answer different questions, and a reader should not take
either alone as the summary.

One further comparison is worth stating plainly because it looks damning in
isolation. On raw \(R^2\) M100 (\num{-0.189}) transfers \emph{worse} than
synthetic data with no relation at all (\texttt{no\_law}, \num{-0.070}). That
is not a contradiction --- M100's ordering correlation is five times
\texttt{no\_law}'s --- but it is a direct demonstration that raw \(R^2\) is the
wrong summary for this question, since a large family-specific level offset
punishes it more than having no relation does.

\subsection{Predictability transfers ordinally, and only partly in level}
\label{sec:linking}

Individually the measures correlate with skill. Cluster-robust over \num{261}
metric clusters the strongest is the gzip compression ratio at \(r=-0.42\)
(95\% CI \([-0.52,-0.31]\)), ahead of repetition fraction \(+0.37\)
(\([0.26,0.49]\)), permutation entropy \(-0.33\) and changepoint rate
\(-0.30\). That a general-purpose compressor outranks every hand-designed
entropy measure is worth noting on its own.

Jointly the same pattern holds and is stronger. The three highest permutation
importances in the fitted model are all information-theoretic and none is
hand-designed for time series: the compression ratio (\num{0.110}), the
Lempel-Ziv entropy rate (\num{0.048}), and the Fano ceiling derived from it
(\num{0.037}). Repetition fraction, which an earlier draft of this work reported
as the dominant predictor, falls to \num{0.019} once they are present --- its
signal was largely theirs. Adding the three lifts every transfer statistic in
Table~\ref{tab:transfer}: median \(R^2\) from \(+0.05\) to \(+0.14\), families
with positive within-family ordering from \num{14} of \num{15} to \emph{all}
\num{15}, and median MAE-skill from \num{0.19} to \num{0.21}.

Jointly, the raw picture looks poor: holding out each of \num{15} families in
turn gives a mean out-of-sample \(R^2\) of \(-0.17\), a median of \(+0.14\),
positive for \num{10} of the \num{15}, ranging from \(+0.42\) to \(-2.56\)
(monotone error counters). Reported alone that number is
misleading in three ways, and Table~\ref{tab:transfer} shows why.

\emph{\(R^2\) is scored against an oracle, not against a usable baseline.}
\(R^2\) is \(1-\mathrm{SSE}/\mathrm{SST}\), and \(\mathrm{SST}\) is taken about
the mean of the \emph{held-out} family. Its implicit comparator is therefore a
predictor that already knows the average skill of the family being predicted,
which is precisely the quantity unavailable to a model that has never seen it.
\(R^2<0\) means \emph{worse than that oracle}; it does not mean \emph{no better
than nothing}. The distinction is not academic here, because the denominator
degenerates whenever a family's series have similar skill. Alarm states are the
clearest case: their skills have a standard deviation of \num{0.036}, the model
predicts every held-out series to a mean absolute error of \num{0.036} against
\num{0.600} for the best family-blind predictor --- a \num{94}\% reduction ---
and \(R^2\) records \(-2.05\).

Scoring the same folds the way this paper scores forecasts everywhere else, as
a fractional reduction in absolute error against a baseline
(\(1-\mathrm{MAE}_{\text{model}}/\mathrm{MAE}_{\text{baseline}}\), the same
construction as MASE), the picture inverts: \num{13} of \num{15} families beat
the family-blind predictor, by a median \num{21}\% reduction in error. The one
substantive failure is monotone error counters, at \num{68}\% \emph{worse} than
baseline. The synthetic control of \S\ref{sec:protocol} makes the point sharpest:
the arm with skill permuted across series, containing no relation whatever,
scores \(R^2=-0.07\) --- \emph{better} than M100's \(-0.17\) --- while its
MAE-skill is \(+0.002\) against M100's \(+0.21\) and its within-family Spearman
\(+0.10\) against \(+0.52\). A statistic that ranks genuine structure below
noise is not measuring the question being asked, and we report it here only
because it is what a reader would expect to see.

\emph{Ordering transfers; level does not.} Within a held-out family the model
ranks series correctly --- median Spearman \(+0.52\), positive in all
\num{15} families and significant at \(p<0.05\) for \num{10} --- while
placing that ranking at the wrong absolute level. \(R^2\) penalises the level
error; Spearman does not. Removing the level bias moves the median \(R^2\) from
\(+0.14\) to \(+0.15\) and makes \num{11} of \num{15} families positive. The
practical reading is that the profile can say which signals in an unfamiliar
family will be harder than which, but not how hard any of them will be.

\emph{The failure is specific to family boundaries.} Running the identical
protocol with \emph{random} groups of matched size gives a median \(R^2\) of
\(+0.35\) across draws (2.5--97.5\% range \([0.22, 0.38]\), mean \(+0.32\))
with every group positive,
against \(-0.19\) and \num{60}\% for the real families; over \num{100} draws
no random grouping was as poor as the observed one (\(p<0.001\)). Arbitrary
grouping therefore transfers easily, so the protocol is capable of transfer and
the family axis carries real structure. Without this control the raw \(R^2\)
would be equally consistent with a protocol that simply cannot generalise at
all.

\begin{table}[t]
\centering
\caption{Transfer to a held-out family, under three criteria. Against a
baseline a model could actually have, transfer succeeds for \num{13} of
\num{15} families; \(R^2\) disagrees because its comparator is an oracle that
already knows the held-out family's mean skill.}
\label{tab:transfer}
\small
\begin{tabular}{lrrr}
\toprule
 & mean & median & positive \\
\midrule
\multicolumn{4}{l}{\emph{against a family-blind baseline}}\\
MAE-skill                   & \num{+0.218} & \num{+0.213} & 13/15 \\
within-family Spearman      & \num{+0.450} & \num{+0.523} & 10/15 at \(p<0.05\) \\
\midrule
\multicolumn{4}{l}{\emph{against an oracle that knows the family's mean}}\\
raw \(R^2\)                 & \num{-0.167} & \num{+0.141} & 10/15 \\
level-corrected \(R^2\)     & \num{-0.012} & \num{+0.154} & 11/15 \\
\midrule
random-group null (\(R^2\)) & \num{+0.319} & \num{+0.329} & 15/15 \\
\bottomrule
\end{tabular}
\end{table}

\begin{table}[t]
\centering
\caption{Ablations. Removing the entropy family \emph{improves} transfer;
removing the instrument covariates costs almost nothing; using them alone
destroys it.}
\label{tab:ablation}
\small
\begin{tabular}{lrr}
\toprule
Feature set & \# & mean \(R^2\) \\
\midrule
without entropy measures   & 23 & \num{-0.142} \\
full profile               & 27 & \num{-0.167} \\
signal measures only       & 23 & \num{-0.204} \\
instrument covariates only & 4  & \num{-14.999} \\
\midrule
predicting the global mean & --- & \num{-18.352} \\
\bottomrule
\end{tabular}
\end{table}


The result is not an artefact of instrumentation. Using the instrument covariates
alone collapses the model (\(-15.0\) against \(-0.17\) for the full profile),
while removing them costs almost nothing (\(-0.20\)) --- see
Table~\ref{tab:ablation}. Whatever the
measures capture, it is not how the data was sampled --- which is the
conditioning test \oi{1} called for.

\subsection{How much labelling closes the level gap}
\label{sec:fewshot}

Leave-one-family-out asks the hardest question available: predict a family
having seen nothing from it. An operator meeting a new signal family would not
refuse to measure a few of its series first. We therefore sweep the number of
labelled series \(k\) drawn at random from the held-out family into the training
set, averaging over \num{15} draws per family and scoring only on the series
that were \emph{not} drawn, so nothing is ever evaluated by a model that saw it.

\begin{table}[t]
\centering
\caption{Transfer as a function of how many series from the held-out family are
labelled first. The level gap is closable, but by steady accumulation rather
than a few calibration points.}
\label{tab:fewshot}
\small
\begin{tabular}{rrr}
\toprule
labelled series \(k\) & MAE-skill & mean \(R^2\) \\
\midrule
0  & \num{0.198} & \num{-0.212} \\
1  & \num{0.208} & \num{-0.170} \\
2  & \num{0.220} & \num{-0.142} \\
5  & \num{0.229} & \num{-0.090} \\
10 & \num{0.254} & \num{-0.015} \\
20 & \num{0.289} & \num{+0.061} \\
\bottomrule
\end{tabular}
\end{table}

Table~\ref{tab:fewshot} shows the gap is real and closable. Twenty labelled
series move the mean \(R^2\) from \(-0.21\) to \(+0.06\) --- past the oracle
threshold the zero-shot result fails --- and raise MAE-skill from \num{0.198} to
\num{0.289}. But the curve is close to linear over this range rather than
saturating early: one or two calibration series buy very little, and there is no
knee after which further labelling stops helping. So the missing level
information is distributed across the family rather than concentrated in a
representative example, which is consistent with the level failure not being an
extrapolation artefact (\S\ref{sec:protocol}).

We note one honest qualification. The number of families with positive mean
\(R^2\) does not rise with \(k\) (nine at \(k=0\), nine at \(k=10\)); the gain
is the poorly-transferring families becoming less poor, not additional families
crossing zero. Monotone error counters remain negative at every \(k\) tested.

\section{Threats to validity}
\label{sec:threats}

The full register is maintained in \texttt{OPEN\_ISSUES.md} and each item is
cited from the text above. The load-bearing ones are:

\begin{description}
  \item[\oi{1}] Ganglia's sampling irregularity confounds instrumentation with
        signal physics. Sampling regularity is therefore carried as a covariate
        and reported separately.
  \item[\oi{2}] The monitoring system changes over the campaign.
  \item[\oi{6}] No neural tier has been run; no claim of the form ``not even
        deep models help'' is licensed.
  \item[\oi{7}] The error floor is an estimate with a known bias direction, not
        an information-theoretic bound.
  \item[\oi{11}] Permutation entropy is biased downward by quantisation, and
        quantisation varies by family, so entropy is interpreted only alongside
        the measured quantisation step.
\end{description}

\section{Conclusion}

We set out to show that the difficulty of forecasting HPC telemetry is an
intrinsic, measurable property of the signals. Half of that survives. Difficulty
is not a modelling deficiency: across four tiers spanning naive baselines to
neural networks, more capacity produced less skill, and beyond three hours
nothing beat repeating the last value.

The other half holds, in a more specific sense than we set out to show. The
measures transfer to an unseen family in \emph{ordering} --- median
within-family Spearman \(+0.52\), positive in \emph{all} \num{15} families
and significant for \num{10} --- and in \emph{level} they beat the best
family-blind predictor in \num{13} of \num{15} families by a median \num{21}\%,
while falling well short of an oracle that knows the family's mean. Two controls
confirm this is a property of the data and not of the method: arbitrary groups
transfer easily (\(R^2 = +0.35\), \(p<0.001\)), and a synthetic portable law is
recovered by the identical pipeline at \(R^2 = +0.25\), matching an oracle
handed the true generating parameter. So there is a portable component of
forecast difficulty; it is strongest ordinally and real but partial in level.

A methodological point travels with that result, and we think it generalises
beyond this dataset. Reported as a raw \(R^2\), our transfer looks worse than
the synthetic arm in which skill was permuted at random and no relation exists
at all (\(-0.17\) against \(-0.07\)). \(R^2\)'s comparator is an oracle that
already knows the mean of the group being predicted, and its denominator
degenerates on homogeneous groups --- our alarm-state family is predicted to a
mean absolute error of \num{0.036} against a baseline of \num{0.600} and scores
\(R^2=-2.05\). Grouped cross-validation summarised by \(R^2\) will
systematically understate transfer wherever groups are internally homogeneous,
which in telemetry they very often are.

What blocks the level from transferring, we do not know. The natural
explanation --- that families occupy disjoint bands of a shared law, so a
held-out family requires extrapolation --- is ruled out by the same controls:
in a synthetic disjoint world the ordering collapses along with the level, and
M100's ordering does not. We leave the mechanism open rather than assert the
one the evidence excludes.

That distinction is the practical result. An operator facing an unfamiliar
signal family can use a model-free profile to decide which of its signals to
attempt first, and can estimate how hard they will be substantially better than
by guessing --- but not as well as someone who has already worked with that
family. The remaining gap is closable by labelling, but not cheaply: twenty
series measured from the new family move the mean \(R^2\) from \(-0.21\) to
\(+0.06\), and the curve is near-linear over that range, so there is no small
calibration set that suffices (\S\ref{sec:fewshot}).

For operators the practical reading is narrower than we hoped but still useful:
on this data, at these horizons, a persistence forecast is a strong default ---
significantly better than every heavier tier we tested, at every horizon --- and
the burden of proof lies with any claim that a larger model improves on it.

\section*{Disclosure}

The author is a co-author of the M100 ExaData dataset paper~\citep{borghesi2023m100}
analysed here. No part of this study was performed as part of that data-collection
campaign, and the analysis uses only the public release.

\section*{Data and code availability}

All code, the versioned literature query set, the screening table, and every
derived table are at
\url{https://github.com/MSKazemi/hpc-signal-forecastability}. The underlying
measurements are the public M100 ExaData release~\citep{borghesi2023m100}.
Results reproduce from the committed derived tables without downloading raw
data.

\bibliographystyle{plainnat}
\bibliography{refs}

\end{document}
